% Part: lambda-calculus
% Chapter: introduction
% Section: term-revisited

\documentclass[../../../include/open-logic-section]{subfiles}

\begin{document}

\olfileid[pt]{lam}{syn}{tr}
\olsection{Termos como classes de equivalência~$\alpha$}

Daqui em diante, consideraremos os termos a menos de equivalência~$\alpha$.
Isso significa que, quando escrevemos um termo, queremos dizer a classe de
equivalência~$\alpha$ à qual ele pertence. Por exemplo, escrevemos
$\lambd[a][\lambd[b][a c]]$ para o conjunto de todos os termos
$\alpha$-equivalentes a ele, como $\lambd[a][\lambd[b][a c]]$,
$\lambd[b][\lambd[a][b c]]$ etc.

Além disso, embora nas seções anteriores letras como $N, Q$ fossem usadas
para denotar um termo, daqui em diante as usaremos para denotar uma classe, e
serão essas classes, em vez dos termos, os objetos de estudo do que segue.
Letras como $x, y$ continuam a denotar uma variável.

Também adotamos a notação $\rep{M}$ para denotar um !!{element} arbitrário da
classe~$M$, e $\rep{M}[0], \rep{M}[1], etc. $ caso precisemos de mais de um.

Reutilizamos as notações dos termos para simplificar nossa formulação. Temos
a seguinte definição sobre classes:
\begin{defn}
  \begin{enumerate}
  \item $\lambd[x][N]$ é definida como a classe que contém
    $\lambd[x][\rep{N}]$.
  \item $PQ$ é definida como a classe que contém $\rep{P}\rep{Q}$.
  \end{enumerate}
\end{defn}

Não é difícil ver que estão bem definidas, pois a conversão~$\alpha$ é
compatível.

\begin{defn} \ollabel{def:fv}
  As \emph{variáveis livres} de uma classe de equivalência~$\alpha$~$M$, ou
  $FV(M)$, são definidas como $FV(\rep{M})$.
\end{defn}

Isso está bem definido, pois $FV(\rep{M}[0]) = FV(\rep{M}[1])$, como mostrado
em \olref[alp]{thm:fv}.

Também reutilizamos a notação para substituição em classes:
\begin{defn} \ollabel{def:sub}
  A \emph{substituição} de $R$ no lugar de $y$ em $M$, ou $\Subst{M}{R}{y}$, é
  definida como $\Subst{\rep{M}}{\rep{R}}{y}$, para quaisquer $\rep{M}$ e
  $\rep{R}$ que tornem a substituição definida.
\end{defn}

Isso também está bem definido, como mostrado em \olref[alp]{cor:sub}.

Observe como essa definição simplifica significativamente nosso raciocínio.
Por exemplo:
\begin{align}
  \Subst{\lambd[x][x]}{y}{x} & =\ollabel{eq:1}\\
  &= \Subst{\lambd[z][z]}{y}{x} \ollabel{eq:2}\\
  &= \lambd[z][\Subst{z}{y}{x}] \\
  &= \lambd[z][z]
\end{align}

\olref{eq:1} não está definida se ainda a considerarmos uma substituição
sobre termos; mas, como mencionado anteriormente, agora a consideramos uma
substituição sobre classes, razão pela qual \olref{eq:2} pode ocorrer: podemos
substituir $\lambd[x][x]$ por $\lambd[z][z]$ porque pertencem à mesma classe.

Pela mesma razão, daqui em diante suporemos que os representantes escolhidos
sempre satisfaçam as condições necessárias à substituição. Por exemplo, ao
vermos $\Subst{\lambd[x][N]}{R}{y}$, suporemos que o representante
$\lambd[x][N]$ foi escolhido de modo que $x \neq y$ e $x \notin FV(R)$.

Como é um pouco estranho chamar $\lambd[x][x]$ de «classe», daqui em diante
vamos chamá-los de termos~$\Lambda$ (ou simplesmente «termos» no restante
desta parte), para distingui-los dos termos~$\lambda$ com os quais estamos
familiarizados.

\begin{editorial}
  Ainda não podemos nos despedir dos termos: toda a definição dos
  termos~$\Lambda$ baseia-se nos termos~$\lambda$, e não fornecemos um método
  para definir funções sobre termos~$\Lambda$. Isso significa que todas essas
  funções precisam ser definidas primeiro sobre termos~$\lambda$ e depois
  «projetadas» para os termos~$\Lambda$, como fizemos com as substituições.
  Contudo, supomos que o leitor possa compreender intuitivamente como definir
  funções sobre termos~$\Lambda$.
\end{editorial}
\end{document}
